\begin{proof}
\oldpage[IV-2]{178}
(i) Si $f$ es radicial en el punto $y$ y $g$ es radicial en el
punto $z$, entonces $g^{-1}(z)$ se reduce a $y$, y
$f^{-1}(g^{-1}(z))=f^{-1}(y)$ es vacío o se reduce a un único
punto $x$. Además, $\boldsymbol{k}(x)$ es una extensión radicial de $\boldsymbol{k}(y)$,
y $\boldsymbol{k}(y)$ una extensión radicial de $\boldsymbol{k}(z)$; por tanto, $\boldsymbol{k}(x)$ es
una extensión radicial de $\boldsymbol{k}(z)$. Recíprocamente, supongamos
que $f$ es sobreyectivo. Si $g\circ f$ es radicial en el punto
$z=g(y)$, entonces $g^{-1}(z)$ se reduce al único punto $y$;
de lo contrario, $f^{-1}(g^{-1}(z))$ contendría al menos dos
puntos distintos. Además,
$f^{-1}(y)=f^{-1}(g^{-1}(z))$ se reduce a un único punto $x$.
Por hipótesis, $\boldsymbol{k}(z)\subset \boldsymbol{k}(y)\subset \boldsymbol{k}(x)$ y $\boldsymbol{k}(x)$ es
radicial sobre $\boldsymbol{k}(z)$; de ahí que $\boldsymbol{k}(y)$ sea radicial sobre
$\boldsymbol{k}(z)$ y $\boldsymbol{k}(x)$ lo sea sobre $\boldsymbol{k}(y)$.

(ii) Sean $g:f^{-1}(y)\to\operatorname{Spec}(\boldsymbol{k}(y))$ y
$g':f'^{-1}(y')\to\operatorname{Spec}(\boldsymbol{k}(y'))$ los morfismos
obtenidos de $f$ y $f'$, respectivamente, por cambio de base.
Como $f'^{-1}(y')=f^{-1}(y)\otimes_{\mathbf{k}(y)}\boldsymbol{k}(y')$ (\textbf I, 3.6.4),
$g'$ se obtiene de $g$ por cambio de base, y la afirmación se
reduce a (2.6.1, (v)).
\end{proof}

Sea ahora $X$ un preesquema \emph{reducido} con un número \emph{finito} de
componentes irreducibles, y sea $X'$ su \emph{normalizado} (\textbf{II}, 6.3.8);
se sabe \emph{(loc. cit.)} que el morfismo canónico $f:X'\to X$ es
sobreyectivo. La definición (6.15.1) muestra entonces que, para
que $X$ sea geométricamente unirramoso en un punto $x$, es
necesario y suficiente que $f$ sea \emph{radicial en ese punto}; por
tanto, $X$ es geométricamente unirramoso si y solo si $f$ es
\emph{radicial}.
